% Options for packages loaded elsewhere
\PassOptionsToPackage{unicode}{hyperref}
\PassOptionsToPackage{hyphens}{url}
\documentclass[
]{article}
\usepackage{xcolor}
\usepackage[margin=1in]{geometry}
\usepackage{amsmath,amssymb}
\setcounter{secnumdepth}{-\maxdimen} % remove section numbering
\usepackage{iftex}
\ifPDFTeX
  \usepackage[T1]{fontenc}
  \usepackage[utf8]{inputenc}
  \usepackage{textcomp} % provide euro and other symbols
\else % if luatex or xetex
  \usepackage{unicode-math} % this also loads fontspec
  \defaultfontfeatures{Scale=MatchLowercase}
  \defaultfontfeatures[\rmfamily]{Ligatures=TeX,Scale=1}
\fi
\usepackage{lmodern}
\ifPDFTeX\else
  % xetex/luatex font selection
\fi
% Use upquote if available, for straight quotes in verbatim environments
\IfFileExists{upquote.sty}{\usepackage{upquote}}{}
\IfFileExists{microtype.sty}{% use microtype if available
  \usepackage[]{microtype}
  \UseMicrotypeSet[protrusion]{basicmath} % disable protrusion for tt fonts
}{}
\makeatletter
\@ifundefined{KOMAClassName}{% if non-KOMA class
  \IfFileExists{parskip.sty}{%
    \usepackage{parskip}
  }{% else
    \setlength{\parindent}{0pt}
    \setlength{\parskip}{6pt plus 2pt minus 1pt}}
}{% if KOMA class
  \KOMAoptions{parskip=half}}
\makeatother
\usepackage{color}
\usepackage{fancyvrb}
\newcommand{\VerbBar}{|}
\newcommand{\VERB}{\Verb[commandchars=\\\{\}]}
\DefineVerbatimEnvironment{Highlighting}{Verbatim}{commandchars=\\\{\}}
% Add ',fontsize=\small' for more characters per line
\usepackage{framed}
\definecolor{shadecolor}{RGB}{248,248,248}
\newenvironment{Shaded}{\begin{snugshade}}{\end{snugshade}}
\newcommand{\AlertTok}[1]{\textcolor[rgb]{0.94,0.16,0.16}{#1}}
\newcommand{\AnnotationTok}[1]{\textcolor[rgb]{0.56,0.35,0.01}{\textbf{\textit{#1}}}}
\newcommand{\AttributeTok}[1]{\textcolor[rgb]{0.13,0.29,0.53}{#1}}
\newcommand{\BaseNTok}[1]{\textcolor[rgb]{0.00,0.00,0.81}{#1}}
\newcommand{\BuiltInTok}[1]{#1}
\newcommand{\CharTok}[1]{\textcolor[rgb]{0.31,0.60,0.02}{#1}}
\newcommand{\CommentTok}[1]{\textcolor[rgb]{0.56,0.35,0.01}{\textit{#1}}}
\newcommand{\CommentVarTok}[1]{\textcolor[rgb]{0.56,0.35,0.01}{\textbf{\textit{#1}}}}
\newcommand{\ConstantTok}[1]{\textcolor[rgb]{0.56,0.35,0.01}{#1}}
\newcommand{\ControlFlowTok}[1]{\textcolor[rgb]{0.13,0.29,0.53}{\textbf{#1}}}
\newcommand{\DataTypeTok}[1]{\textcolor[rgb]{0.13,0.29,0.53}{#1}}
\newcommand{\DecValTok}[1]{\textcolor[rgb]{0.00,0.00,0.81}{#1}}
\newcommand{\DocumentationTok}[1]{\textcolor[rgb]{0.56,0.35,0.01}{\textbf{\textit{#1}}}}
\newcommand{\ErrorTok}[1]{\textcolor[rgb]{0.64,0.00,0.00}{\textbf{#1}}}
\newcommand{\ExtensionTok}[1]{#1}
\newcommand{\FloatTok}[1]{\textcolor[rgb]{0.00,0.00,0.81}{#1}}
\newcommand{\FunctionTok}[1]{\textcolor[rgb]{0.13,0.29,0.53}{\textbf{#1}}}
\newcommand{\ImportTok}[1]{#1}
\newcommand{\InformationTok}[1]{\textcolor[rgb]{0.56,0.35,0.01}{\textbf{\textit{#1}}}}
\newcommand{\KeywordTok}[1]{\textcolor[rgb]{0.13,0.29,0.53}{\textbf{#1}}}
\newcommand{\NormalTok}[1]{#1}
\newcommand{\OperatorTok}[1]{\textcolor[rgb]{0.81,0.36,0.00}{\textbf{#1}}}
\newcommand{\OtherTok}[1]{\textcolor[rgb]{0.56,0.35,0.01}{#1}}
\newcommand{\PreprocessorTok}[1]{\textcolor[rgb]{0.56,0.35,0.01}{\textit{#1}}}
\newcommand{\RegionMarkerTok}[1]{#1}
\newcommand{\SpecialCharTok}[1]{\textcolor[rgb]{0.81,0.36,0.00}{\textbf{#1}}}
\newcommand{\SpecialStringTok}[1]{\textcolor[rgb]{0.31,0.60,0.02}{#1}}
\newcommand{\StringTok}[1]{\textcolor[rgb]{0.31,0.60,0.02}{#1}}
\newcommand{\VariableTok}[1]{\textcolor[rgb]{0.00,0.00,0.00}{#1}}
\newcommand{\VerbatimStringTok}[1]{\textcolor[rgb]{0.31,0.60,0.02}{#1}}
\newcommand{\WarningTok}[1]{\textcolor[rgb]{0.56,0.35,0.01}{\textbf{\textit{#1}}}}
\usepackage{longtable,booktabs,array}
\newcounter{none} % for unnumbered tables
\usepackage{calc} % for calculating minipage widths
% Correct order of tables after \paragraph or \subparagraph
\usepackage{etoolbox}
\makeatletter
\patchcmd\longtable{\par}{\if@noskipsec\mbox{}\fi\par}{}{}
\makeatother
% Allow footnotes in longtable head/foot
\IfFileExists{footnotehyper.sty}{\usepackage{footnotehyper}}{\usepackage{footnote}}
\makesavenoteenv{longtable}
\usepackage{graphicx}
\makeatletter
\newsavebox\pandoc@box
\newcommand*\pandocbounded[1]{% scales image to fit in text height/width
  \sbox\pandoc@box{#1}%
  \Gscale@div\@tempa{\textheight}{\dimexpr\ht\pandoc@box+\dp\pandoc@box\relax}%
  \Gscale@div\@tempb{\linewidth}{\wd\pandoc@box}%
  \ifdim\@tempb\p@<\@tempa\p@\let\@tempa\@tempb\fi% select the smaller of both
  \ifdim\@tempa\p@<\p@\scalebox{\@tempa}{\usebox\pandoc@box}%
  \else\usebox{\pandoc@box}%
  \fi%
}
% Set default figure placement to htbp
\def\fps@figure{htbp}
\makeatother
\setlength{\emergencystretch}{3em} % prevent overfull lines
\providecommand{\tightlist}{%
  \setlength{\itemsep}{0pt}\setlength{\parskip}{0pt}}
\usepackage{bookmark}
\IfFileExists{xurl.sty}{\usepackage{xurl}}{} % add URL line breaks if available
\urlstyle{same}
\hypersetup{
  hidelinks,
  pdfcreator={LaTeX via pandoc}}

\author{}
\date{\vspace{-2.5em}}

\begin{document}

{\def\LTcaptype{none} % do not increment counter
\begin{longtable}[]{@{}ccc@{}}
\toprule\noalign{}
Version & Status & Date \\
\midrule\noalign{}
\endhead
\bottomrule\noalign{}
\endlastfoot
0.0.1 & draft & 2026-08-18 \\
1.0.0 & final & 2026-08-21 \\
1.1.0 & update & 2026-08-27 \\
1.2.0 & update & 2026-08-31 \\
\end{longtable}
}

\section{Eurostat Metabase Structure
Monitor}\label{eurostat-metabase-structure-monitor}

Technical documentation. Covers what the system does, why it is built
the way it is, the data model, the processing pipeline, the subscription
mechanism, operational details, and known limitations.

The monitor draws on \textbf{two independent Eurostat sources}: the
\textbf{metabase} (the full dimensional structure of every dataset ---
Sections 2, 4.2, 5, 6) and the \textbf{TOC} (the navigation-tree
hierarchy of folders and datasets --- Sections 2b, 4.5, 5b). The
metabase drives structural-change detection; the TOC lets users
subscribe to whole folders rather than listing individual datasets
(Section 7.3).

\begin{center}\rule{0.5\linewidth}{0.5pt}\end{center}

\subsection{1. Purpose and problem}\label{purpose-and-problem}

Eurostat datasets change structure without notice: dimensions are
renamed or added, code lists gain or lose members, whole datasets are
renamed or withdrawn. These changes ship silently with Eurostat's
twice-daily refresh (11:00 and 23:00 CET). There is \textbf{no
forward-looking schedule} of structural changes anywhere in Eurostat's
dissemination system, and the database keeps only the latest version of
each dataset --- there is no versioning or history on their side.

Downstream, this breaks things:

\begin{itemize}
\tightlist
\item
  Colleagues' Power Query links that hard-code dimension names and level
  codes in the API URL fail the moment a dimension is renamed or added.
\item
  Automated R ingestion pipelines that assume a fixed structure break
  or, worse, silently mis-ingest.
\end{itemize}

The monitor's job is to \textbf{detect structural changes as soon as
they publish}, record them as durable history (which Eurostat itself
does not keep), and \textbf{alert the specific people who depend on the
affected datasets} --- before they discover the breakage the hard way.

The system does not prevent breakage (it cannot; a renamed dimension
breaks any query that references it). It makes breakage \textbf{visible
early and legible}, and routes each change to the people who care about
it.

\begin{center}\rule{0.5\linewidth}{0.5pt}\end{center}

\subsection{2. Data source}\label{data-source}

The single source is the Eurostat \textbf{metabase} file:

\begin{verbatim}
https://ec.europa.eu/eurostat/api/dissemination/catalogue/metabase.txt.gz
\end{verbatim}

It is a gzipped, tab-separated file giving the full structure of every
disseminated dataset, one row per
\texttt{(dataset,\ dimension,\ position)}:

\begin{verbatim}
aact_ali01   geo    EU
aact_ali01   geo    EU27_2020
aact_ali01   itm_newa   40000
aact_ali01   time   2022
...
\end{verbatim}

At the time of writing it is \textasciitilde1.34M rows across
\textasciitilde7,600 datasets. Rows are ordered by dataset, then
dimension, then the code list's protocol order --- this row order is
itself information (it is the intended order of the code list) and is
captured on ingest.

\subsubsection{Source constraints that shaped the
design}\label{source-constraints-that-shaped-the-design}

\begin{itemize}
\tightlist
\item
  \textbf{No \texttt{Last-Modified} header.} The endpoint returns a
  chunked, generated response, so there is no file modification
  timestamp to use for change detection or as a publication time. This
  is confirmed behaviour, not a bug in the fetch. Consequently the
  observation timestamp we record is \emph{our} fetch time, not
  Eurostat's publication time.
\item
  \textbf{\texttt{NA} collides with Namibia.} Namibia's ISO-2 code is
  the literal string \texttt{NA}, which \texttt{data.table::fread}
  converts to a missing value by default. This must be suppressed on
  read (\texttt{na.strings\ =\ NULL}), or Namibia silently vanishes from
  every dataset that has it as a position, corrupting the structure
  record.
\item
  \textbf{The metabase carries codes, not labels.} It gives dimension
  \emph{codes} and position \emph{codes}, never their human-readable
  titles. This is why some kinds of rename cannot be detected (see
  Limitations).
\end{itemize}

\begin{center}\rule{0.5\linewidth}{0.5pt}\end{center}

\subsection{2b. Second data source: the TOC (navigation
tree)}\label{b.-second-data-source-the-toc-navigation-tree}

The second source is the Eurostat \textbf{table of contents} (navigation
tree) XML:

\begin{verbatim}
https://ec.europa.eu/eurostat/api/dissemination/catalogue/toc/xml
\end{verbatim}

It is the folder hierarchy the online data browser shows: nested
\texttt{branch} (folder) elements containing \texttt{leaf} (dataset or
table) elements, each with a stable \textbf{code} and
English/French/German titles. Both branches and leaves may contain
children (a dataset leaf can nest curated \texttt{table} leaves), so it
is a genuine tree of arbitrary depth. At the time of writing it parses
to \textasciitilde9,400 distinct nodes and \textasciitilde12,200 edges.

The TOC exists in the system for one user-facing feature:
\textbf{subscribing to a whole folder} instead of listing every dataset
under it (Section 7.3). It is diffed and snapshotted on the same SCD2
principles as the metabase, but the two are entirely independent sources
with independent snapshots and cadences.

\subsubsection{Source characteristics that shaped the
design}\label{source-characteristics-that-shaped-the-design}

\begin{itemize}
\tightlist
\item
  \textbf{Folders have codes.} Unlike the flat tsv TOC that
  \texttt{eurostat::get\_eurostat\_toc()} parses, the XML carries a
  \texttt{\textless{}code\textgreater{}} on every folder node, so folder
  identity is a stable code, not a mutable title. This is why the ingest
  is a fresh XML parse rather than a reuse of the older tsv-based
  hierarchy code.
\item
  \textbf{The XML carries a real \texttt{creationDate}} (on the
  \texttt{\textless{}tree\textgreater{}} root, format
  \texttt{\%Y\%m\%dT\%H\%M}). Unlike the metabase (no
  \texttt{Last-Modified}), this gives a genuine publication timestamp
  --- so the TOC ingest can run a \emph{monotonicity guard} (refuse a
  fetched tree stamped older than the latest stored snapshot).
\item
  \textbf{Nodes are reached by multiple paths.} \textasciitilde2,000
  codes appear under more than one parent (the same dataset filed in
  several folders). The parser emits one node per code (deduplicated on
  code/type/title) but one edge per parent→child pair, so
  multi-parenting is represented natively as multiple edges.
\item
  \textbf{English titles only} are stored; a node with no English title
  stores \texttt{NA} rather than falling back to another language.
\item
  \textbf{The TOC XML is an incomplete listing relative to the browser.}
  A handful of tables that exist in the metabase (and in the browser)
  are simply absent from the XML file --- confirmed by grepping the raw
  source, not a parser fault. See Limitations.
\end{itemize}

\begin{center}\rule{0.5\linewidth}{0.5pt}\end{center}

\subsection{3. Architecture overview}\label{architecture-overview}

The system separates into \textbf{reusable database components} (which
belong in the R package suite) and an \textbf{application} (the
orchestration and alerting, which is a standalone scheduled script).
This separation is deliberate: the DB-generic pieces are reused and
warrant roxygen documentation and tests; the orchestration is one
specific job and is allowed to be a script.

{\def\LTcaptype{none} % do not increment counter
\begin{longtable}[]{@{}
  >{\raggedright\arraybackslash}p{(\linewidth - 4\tabcolsep) * \real{0.3333}}
  >{\raggedright\arraybackslash}p{(\linewidth - 4\tabcolsep) * \real{0.3333}}
  >{\raggedright\arraybackslash}p{(\linewidth - 4\tabcolsep) * \real{0.3333}}@{}}
\toprule\noalign{}
\begin{minipage}[b]{\linewidth}\raggedright
Concern
\end{minipage} & \begin{minipage}[b]{\linewidth}\raggedright
Belongs in
\end{minipage} & \begin{minipage}[b]{\linewidth}\raggedright
Nature
\end{minipage} \\
\midrule\noalign{}
\endhead
\bottomrule\noalign{}
\endlastfoot
TOC XML fetch + parse (\texttt{fetch\_toc} / \texttt{parse\_toc}) &
EUROSTATfetchR & reusable, tested \\
Change-detection reads (\texttt{sql\_get\_*},
\texttt{resolve\_folder\_datasets}) & UMARaccessR & reusable, tested \\
Generic query executor (\texttt{sql\_function\_call}) & UMARimportR &
reusable \\
Metabase + TOC ingest orchestration & standalone script & application
glue \\
Subscription load/validate, filtering, email & standalone script &
application glue \\
\end{longtable}
}

The DB read functions are Postgres functions (their SQL tracked in
\texttt{EUROSTATfetchR/inst/sql/} alongside the schema DDL, with
\texttt{BEGIN…ROLLBACK} SQL tests), wrapped by thin \texttt{sql\_get\_*}
R functions in UMARaccessR that call
\texttt{UMARimportR::sql\_function\_call}. The read wrappers are grouped
by \emph{what they do} (reads → UMARaccessR), not by which schema they
read --- consistent with the \texttt{platform}-schema read wrappers
living there too.

The application script wires these together and is what Windows Task
Scheduler runs. It is not a package because Task Scheduler runs a
script, and because its email/orchestration logic should not couple the
general-purpose DB libraries to this one monitoring use case. The
metabase SCD2 upsert and the TOC ingest both live in the application as
orchestration (multi-statement, branch logic, source- specific) rather
than as reusable primitives --- a deliberate application of the ``only
genuinely reusable things become package functions'' rule.

All database objects live in a dedicated \texttt{eurostat} schema,
separate from the \texttt{platform} schema that holds the main
statistical data model. The \texttt{eurostat} schema DDL (metabase
tables, TOC tables, and all read functions) is tracked in
\texttt{EUROSTATfetchR/inst/sql/}.

\subsubsection{3.1 File and code layout}\label{file-and-code-layout}

The application is a set of standalone R files sourced by a runner, plus
a batch wrapper for Task Scheduler. Nothing in the application is a
package; the package code lives in the three R packages (EUROSTATfetchR,
UMARaccessR, UMARimportR).

\begin{verbatim}
O:\Avtomatizacija\umar-automation-scripts\
  R\
    eurostat-update-metabase.R          <- runner: sources the function files, calls main()
    eurostat-metabase\                  <- the application function files
      ingest.R          -- ingest() (metabase) and ingest_toc()
      subscriptions.R   -- load_subscriptions(), subscribe_bot_to_ingested()
      alerts.R          -- is_breaking(), filter_alerts()
      email.R           -- init_email(), send_failure_email(),
                           format_alert_body(), send_alerts_perperson()
      main.R            -- main() (defines only; does not call itself)
  batch-scripts\
    eurostat-update-metabase.bat        <- Task Scheduler entry point
  data\umar_data_bot_subscriptions\
    eurostat_email_prijave.xlsx         <- the subscription spreadsheet
  logs\
    eurostat_metabase.Rout              <- redirected run log, committed to git
\end{verbatim}

The runner (\texttt{eurostat-update-metabase.R}, outside the function
folder) sources all files in \texttt{eurostat-metabase\textbackslash{}}
and then calls \texttt{main()}. The function files are pure definitions
--- no \texttt{library()}, no \texttt{source()}, no top-level execution
--- so sourcing them has no side effects and their order does not matter
(\texttt{main()} is only evaluated once called). Package functions are
referenced with \texttt{::}, so only the packages need to be installed
in the scheduled account's library, not attached. The batch file runs
the runner via \texttt{Rscript}, redirects output to the log, and
commits the log to git.

\begin{center}\rule{0.5\linewidth}{0.5pt}\end{center}

\subsection{4. Data model}\label{data-model}

The \texttt{eurostat} schema holds two families of tables: the
\textbf{metabase} family (snapshot + SCD2 structure, Sections 4.1--4.4)
and the \textbf{TOC} family (snapshot + SCD2 node/edge tree, Section
4.5), plus the \texttt{subscription} table shared by both. Each family
has its own temporary staging table used only during ingest.

\subsubsection{\texorpdfstring{4.1 \texttt{snapshot} --- one row per
accepted
ingest}{4.1 snapshot --- one row per accepted ingest}}\label{snapshot-one-row-per-accepted-ingest}

\begin{Shaded}
\begin{Highlighting}[]
\KeywordTok{CREATE} \KeywordTok{TABLE}\NormalTok{ eurostat.}\KeywordTok{snapshot}\NormalTok{ (}
\NormalTok{  snapshot\_id bigserial }\KeywordTok{PRIMARY} \KeywordTok{KEY}\NormalTok{,}
\NormalTok{  file\_sha256 text }\KeywordTok{NOT} \KeywordTok{NULL}\NormalTok{,          }\CommentTok{{-}{-} content identity of the ingested file}
\NormalTok{  observed\_at timestamptz }\KeywordTok{NOT} \KeywordTok{NULL} \KeywordTok{DEFAULT}\NormalTok{ now(),}
\NormalTok{  n\_rows      bigint }\KeywordTok{NOT} \KeywordTok{NULL}
\NormalTok{);}
\end{Highlighting}
\end{Shaded}

\begin{itemize}
\tightlist
\item
  \texttt{file\_sha256} is the \textbf{content identity} of the metabase
  file. It is the dedup key: if the newly fetched file's hash matches
  the latest snapshot, there is nothing to do.
\item
  \texttt{file\_sha256} is \textbf{not} unique. A rollback (Eurostat
  reverting to a byte-identical earlier structure) is a real event we
  want to record as a new observation, so the same hash may legitimately
  appear on multiple dates. Uniqueness is on \texttt{snapshot\_id} (the
  observation), not on content.
\item
  \texttt{observed\_at} is fetch time, not publication time (no header
  available).
\end{itemize}

\subsubsection{\texorpdfstring{4.2 \texttt{metabase} --- the structural
history
(SCD2)}{4.2 metabase --- the structural history (SCD2)}}\label{metabase-the-structural-history-scd2}

\begin{Shaded}
\begin{Highlighting}[]
\KeywordTok{CREATE} \KeywordTok{TABLE}\NormalTok{ eurostat.metabase (}
\NormalTok{  dataset text }\KeywordTok{NOT} \KeywordTok{NULL}\NormalTok{,}
\NormalTok{  dim     text }\KeywordTok{NOT} \KeywordTok{NULL}\NormalTok{,}
\NormalTok{  pos     text }\KeywordTok{NOT} \KeywordTok{NULL}\NormalTok{,}
\NormalTok{  ord     }\DataTypeTok{int}  \KeywordTok{NOT} \KeywordTok{NULL}\NormalTok{,              }\CommentTok{{-}{-} protocol order within (dataset, dim)}
\NormalTok{  valid\_from timestamptz }\KeywordTok{NOT} \KeywordTok{NULL}\NormalTok{,}
\NormalTok{  valid\_to   timestamptz,             }\CommentTok{{-}{-} NULL = live}
\NormalTok{  from\_snapshot\_id bigint }\KeywordTok{NOT} \KeywordTok{NULL} \KeywordTok{REFERENCES}\NormalTok{ eurostat.}\KeywordTok{snapshot}\NormalTok{(snapshot\_id),}
\NormalTok{  to\_snapshot\_id   bigint          }\KeywordTok{REFERENCES}\NormalTok{ eurostat.}\KeywordTok{snapshot}\NormalTok{(snapshot\_id)}
\NormalTok{);}

\KeywordTok{CREATE} \KeywordTok{UNIQUE} \KeywordTok{INDEX}\NormalTok{ ux\_mb\_live }\KeywordTok{ON}\NormalTok{ eurostat.metabase (dataset, dim, pos)}
  \KeywordTok{WHERE}\NormalTok{ valid\_to }\KeywordTok{IS} \KeywordTok{NULL}\NormalTok{;}
\KeywordTok{CREATE} \KeywordTok{INDEX}\NormalTok{ ix\_mb\_from }\KeywordTok{ON}\NormalTok{ eurostat.metabase (from\_snapshot\_id);}
\KeywordTok{CREATE} \KeywordTok{INDEX}\NormalTok{ ix\_mb\_to   }\KeywordTok{ON}\NormalTok{ eurostat.metabase (to\_snapshot\_id);}
\end{Highlighting}
\end{Shaded}

This is a \textbf{Type-2 slowly-changing-dimension} table (SCD2). Each
row represents one \texttt{(dataset,\ dim,\ pos)} combination and the
span over which it was live. A live row has \texttt{valid\_to\ IS\ NULL}
/ \texttt{to\_snapshot\_id\ IS\ NULL}. When a combination disappears
from the file, its row is \emph{expired} (both \texttt{valid\_to} and
\texttt{to\_snapshot\_id} set); when a new combination appears, a new
live row is inserted. Persisting combinations are left untouched.

Design decisions worth recording:

\begin{itemize}
\tightlist
\item
  \textbf{Only changed rows are written.} Each ingest expires the rows
  that vanished and inserts the rows that appeared; the vast unchanged
  majority are never rewritten. This keeps table churn tiny (dozens of
  rows on a normal day) and avoids the MVCC-bloat that an ``expire
  everything, reinsert everything'' approach would cause twice daily.
\item
  \textbf{The unique partial index \texttt{ux\_mb\_live} enforces the
  core invariant:} at most one live row per
  \texttt{(dataset,\ dim,\ pos)}. A duplicate means an ingest bug, and
  the constraint turns that from silent corruption into an immediate,
  transaction-rolling-back failure.
\item
  \textbf{\texttt{from\_snapshot\_id} / \texttt{to\_snapshot\_id} are
  the authoritative lineage keys.} They exist so that ``which ingest
  created/expired this row'' is answered by a foreign key, not by
  matching a timestamp. This is not cosmetic --- see 4.4.
\item
  \textbf{\texttt{ord} captures the code list's protocol order} at the
  time a row was inserted. It is used to build an ordered position
  signature for rename detection. It is \textbf{run-local}: it is not
  renumbered when a code is inserted mid-list in a later run, so
  ordinals can develop duplicates within a long-lived dimension. This is
  acceptable because it only ever makes rename detection
  \emph{conservative} (a false negative, never a false positive), and no
  consumer relies on cross-run ordinal continuity.
  Reorder-without-rename is out of scope.
\item
  \textbf{The \texttt{time} dimension is stored but never alerted on.}
  Every dataset's \texttt{time} dimension gains a position every period
  (new month/quarter/year), so it would be the single largest source of
  false alerts. It is kept in the store (it is real structure) and
  filtered out at detection time.
\end{itemize}

\subsubsection{\texorpdfstring{4.3 \texttt{subscription} --- who wants
to hear about
what}{4.3 subscription --- who wants to hear about what}}\label{subscription-who-wants-to-hear-about-what}

\begin{Shaded}
\begin{Highlighting}[]
\KeywordTok{CREATE} \KeywordTok{TABLE}\NormalTok{ eurostat.subscription (}
\NormalTok{  email   text }\KeywordTok{NOT} \KeywordTok{NULL}\NormalTok{,}
\NormalTok{  dataset text }\KeywordTok{NOT} \KeywordTok{NULL}\NormalTok{,              }\CommentTok{{-}{-} lowercased dataset OR folder code, or \textquotesingle{}*\textquotesingle{}}
  \KeywordTok{scope}\NormalTok{   text }\KeywordTok{NOT} \KeywordTok{NULL} \KeywordTok{DEFAULT} \StringTok{\textquotesingle{}breaking\textquotesingle{}} \KeywordTok{CHECK}\NormalTok{ (}\KeywordTok{scope} \KeywordTok{IN}\NormalTok{ (}\StringTok{\textquotesingle{}breaking\textquotesingle{}}\NormalTok{,}\StringTok{\textquotesingle{}all\textquotesingle{}}\NormalTok{)),}
\NormalTok{  kind    text }\KeywordTok{NOT} \KeywordTok{NULL} \KeywordTok{DEFAULT} \StringTok{\textquotesingle{}dataset\textquotesingle{}} \KeywordTok{CHECK}\NormalTok{ (kind }\KeywordTok{IN}\NormalTok{ (}\StringTok{\textquotesingle{}dataset\textquotesingle{}}\NormalTok{,}\StringTok{\textquotesingle{}folder\textquotesingle{}}\NormalTok{)),}
  \KeywordTok{PRIMARY} \KeywordTok{KEY}\NormalTok{ (email, dataset)}
\NormalTok{);}
\end{Highlighting}
\end{Shaded}

\begin{itemize}
\tightlist
\item
  \texttt{email} is the natural key and the delivery address; there is
  no separate person column.
\item
  \texttt{dataset\ =\ \textquotesingle{}*\textquotesingle{}} subscribes
  to all datasets.
\item
  \texttt{scope} filters by change type: \texttt{breaking} (default)
  sends only changes that can break an existing query; \texttt{all}
  sends everything on the subscribed datasets. See 6.2 for the
  definition of ``breaking''.
\item
  \texttt{kind} is \textbf{derived at load time, not user-supplied}: the
  loader checks each code against the live metabase datasets and the
  live TOC folders and stamps \texttt{dataset} or \texttt{folder}
  accordingly (Section 7.2). The user never declares which --- they just
  enter a code and the backend resolves it. \texttt{*} rows are
  nominally \texttt{dataset} (the \texttt{*} special-case in matching
  fires before \texttt{kind} is consulted).
\item
  The \texttt{CHECK} constraints reject any scope or kind value outside
  the allowed set, so a corrupt load fails loudly rather than producing
  a garbage filter.
\item
  Folder and dataset codes are \textbf{disjoint namespaces} (verified:
  no code is both a live folder and a live dataset), so the derived
  \texttt{kind} is unambiguous.
\item
  Rows are populated from \textbf{two sources}: human subscriptions from
  the Excel sheet (Section 7.1) and machine subscriptions for ingested
  pipeline dependencies (Section 7.4). The two are distinguished only by
  email address --- the bot uses a dedicated address, so its
  \texttt{delete-and-replace} never touches a human's rows.
\end{itemize}

\subsubsection{4.4 Why lineage is keyed on snapshot id, not
timestamp}\label{why-lineage-is-keyed-on-snapshot-id-not-timestamp}

An early version of the change-detection query joined rows to their
snapshot by matching \texttt{valid\_from\ =\ snapshot.observed\_at}.
This is fragile: \texttt{timestamptz} stores microsecond precision, and
two values that \emph{display} identically (both shown to the
millisecond) can differ in lower digits, so the equality silently
matches nothing. This produced a real bug where a query returned zero
rows despite visible changes in the table.

The fix, and the standing rule: \textbf{all change-detection keys off
\texttt{from\_snapshot\_id} / \texttt{to\_snapshot\_id} (integers),
never off the timestamps.} Integer equality has no precision to
mismatch. The timestamp columns remain as human-readable convenience
only; the snapshot ids are authoritative.

\subsubsection{4.5 TOC tables --- the navigation tree
(SCD2)}\label{toc-tables-the-navigation-tree-scd2}

The TOC is stored as a snapshot table plus two SCD2 tables --- nodes and
edges --- so that the tree's history is preserved on the same ``keep
what Eurostat won't'' principle as the metabase.

\begin{Shaded}
\begin{Highlighting}[]
\KeywordTok{CREATE} \KeywordTok{TABLE}\NormalTok{ eurostat.toc\_snapshot (}
\NormalTok{  snapshot\_id   bigserial }\KeywordTok{PRIMARY} \KeywordTok{KEY}\NormalTok{,}
\NormalTok{  file\_sha256   text }\KeywordTok{NOT} \KeywordTok{NULL}\NormalTok{,          }\CommentTok{{-}{-} content identity (dedup key)}
\NormalTok{  creation\_date timestamptz,            }\CommentTok{{-}{-} \textless{}tree creationDate\textgreater{}, provenance}
\NormalTok{  observed\_at   timestamptz }\KeywordTok{NOT} \KeywordTok{NULL} \KeywordTok{DEFAULT}\NormalTok{ now(),}
\NormalTok{  n\_nodes       bigint }\KeywordTok{NOT} \KeywordTok{NULL}
\NormalTok{);}

\KeywordTok{CREATE} \KeywordTok{TABLE}\NormalTok{ eurostat.toc\_node (        }\CommentTok{{-}{-} SCD2; one live row per code}
\NormalTok{  code  text }\KeywordTok{NOT} \KeywordTok{NULL}\NormalTok{,}
  \KeywordTok{type}\NormalTok{  text }\KeywordTok{NOT} \KeywordTok{NULL}\NormalTok{,                  }\CommentTok{{-}{-} \textquotesingle{}folder\textquotesingle{} | \textquotesingle{}dataset\textquotesingle{} | \textquotesingle{}table\textquotesingle{}}
\NormalTok{  title text,}
\NormalTok{  valid\_from timestamptz }\KeywordTok{NOT} \KeywordTok{NULL}\NormalTok{,}
\NormalTok{  valid\_to   timestamptz,               }\CommentTok{{-}{-} NULL = live}
\NormalTok{  from\_snapshot\_id bigint }\KeywordTok{NOT} \KeywordTok{NULL} \KeywordTok{REFERENCES}\NormalTok{ eurostat.toc\_snapshot(snapshot\_id),}
\NormalTok{  to\_snapshot\_id   bigint          }\KeywordTok{REFERENCES}\NormalTok{ eurostat.toc\_snapshot(snapshot\_id)}
\NormalTok{);}
\KeywordTok{CREATE} \KeywordTok{UNIQUE} \KeywordTok{INDEX}\NormalTok{ ux\_toc\_node\_live }\KeywordTok{ON}\NormalTok{ eurostat.toc\_node (code) }\KeywordTok{WHERE}\NormalTok{ valid\_to }\KeywordTok{IS} \KeywordTok{NULL}\NormalTok{;}

\KeywordTok{CREATE} \KeywordTok{TABLE}\NormalTok{ eurostat.toc\_edge (        }\CommentTok{{-}{-} SCD2; one live row per parent→child}
\NormalTok{  parent\_code text }\KeywordTok{NOT} \KeywordTok{NULL}\NormalTok{,}
\NormalTok{  child\_code  text }\KeywordTok{NOT} \KeywordTok{NULL}\NormalTok{,}
\NormalTok{  valid\_from timestamptz }\KeywordTok{NOT} \KeywordTok{NULL}\NormalTok{,}
\NormalTok{  valid\_to   timestamptz,               }\CommentTok{{-}{-} NULL = live}
\NormalTok{  from\_snapshot\_id bigint }\KeywordTok{NOT} \KeywordTok{NULL} \KeywordTok{REFERENCES}\NormalTok{ eurostat.toc\_snapshot(snapshot\_id),}
\NormalTok{  to\_snapshot\_id   bigint          }\KeywordTok{REFERENCES}\NormalTok{ eurostat.toc\_snapshot(snapshot\_id)}
\NormalTok{);}
\KeywordTok{CREATE} \KeywordTok{UNIQUE} \KeywordTok{INDEX}\NormalTok{ ux\_toc\_edge\_live }\KeywordTok{ON}\NormalTok{ eurostat.toc\_edge (parent\_code, child\_code) }
  \KeywordTok{WHERE}\NormalTok{ valid\_to }\KeywordTok{IS} \KeywordTok{NULL}\NormalTok{;}
\end{Highlighting}
\end{Shaded}

Design decisions:

\begin{itemize}
\tightlist
\item
  \textbf{Nodes and edges are separate tables} because the hierarchy
  changes independently of the nodes. A dataset moved from folder A to
  folder B is an \emph{edge} change (expire A→X, insert B→X), not a node
  change. Folding the parent into the node row would make a re-parenting
  look like a node change and would not represent a code with multiple
  parents. Separate edges = clean move detection and native
  multi-parenting (one live edge per parent).
\item
  \textbf{Node identity is \texttt{code}; \texttt{title} and
  \texttt{type} are tracked content.} A title change expires the old row
  and inserts a new one --- the same SCD2 treatment the metabase gives
  every column. Because the TOC (unlike the metabase) carries titles,
  this is what makes \emph{rename detection} possible later: a code with
  an expired title and a new live title is a rename, decisively.
\item
  \textbf{\texttt{creation\_date} enables a monotonicity guard} the
  metabase ingest cannot do: a fetched tree stamped strictly older than
  the latest stored snapshot is refused (a stale-mirror serve). Equal
  timestamps pass (a same-minute regeneration with a real change is
  legitimate).
\item
  Lineage keys off \texttt{from\_snapshot\_id} /
  \texttt{to\_snapshot\_id} throughout, same rule as 4.4.
\end{itemize}

Note on snapshot volume: the TOC XML is regenerated by Eurostat on every
fetch with a fresh \texttt{creationDate}, so its \texttt{sha256} differs
each run even when the tree content is identical. The dedup guard
therefore records a new \texttt{toc\_snapshot} most runs, but a run
whose tree is unchanged expires and inserts \textbf{zero} node and edge
rows --- a true no-op snapshot. These accumulate harmlessly
(\textasciitilde2/day); a run's real changes are exactly those where any
of the node/edge added/expired counts is non-zero. This is accepted
rather than fixed (content-hashing the parsed frames would suppress the
no-op snapshots, but the history is only being accumulated, not yet
diffed, so it is not worth it).

\begin{center}\rule{0.5\linewidth}{0.5pt}\end{center}

\subsection{5. Ingest pipeline}\label{ingest-pipeline}

The ingest is a single function that takes an open connection and
returns a status description. It does not know that email or alerting
exist --- that is the orchestrator's concern. It returns one of:
\texttt{unchanged}, \texttt{changed} (with the new
\texttt{snapshot\_id}), or raises an error.

Steps:

\begin{enumerate}
\def\labelenumi{\arabic{enumi}.}
\tightlist
\item
  \textbf{Fetch} the gzipped metabase to a temp file and compute its
  \texttt{sha256}.
\item
  \textbf{Read} with \texttt{na.strings\ =\ NULL} (Namibia) and capture
  protocol order (\texttt{ord}) per \texttt{(dataset,\ dim)}.
\item
  \textbf{Dedup guard.} Compare the fetched \texttt{sha256} to the
  \emph{latest} snapshot only.

  \begin{itemize}
  \tightlist
  \item
    Matches latest → \texttt{unchanged}, exit. Nothing to do.
  \item
    Matches an \emph{older} snapshot → a rollback. Warn, but proceed to
    ingest it as a new observation.
  \item
    No match → a new structure, proceed.
  \end{itemize}
\item
  \textbf{Diff and apply, in one transaction:}

  \begin{itemize}
  \tightlist
  \item
    Insert the new \texttt{snapshot} row first, capturing its id (needed
    as the FK for the rows about to change).
  \item
    Load the fetched file into an unlogged temp staging table;
    \texttt{ANALYZE} it and index it so the planner uses hash joins, not
    nested loops.
  \item
    \textbf{Bootstrap branch} (\texttt{n\_live\ ==\ 0}): the live set is
    empty, so every row is new by definition --- a plain bulk insert, no
    anti-join. This exists because an anti-join against an empty target
    degenerates to a per-row scan and takes minutes; the plain insert
    takes seconds.
  \item
    \textbf{Steady-state branch:} expire live rows absent from the file
    (anti-join staging → metabase), then insert file rows absent from
    the live set (anti-join metabase → staging). Both are sub-second at
    this scale with the FK indexes present.
  \end{itemize}
\item
  \textbf{Return} \texttt{changed} with the new \texttt{snapshot\_id},
  \texttt{n\_rows}, and a \texttt{rolled\_back} flag.
\end{enumerate}

The whole thing is wrapped so that on any exit --- success, error, or
early return --- the connection is disconnected. A \textbf{session-level
advisory lock} taken at the start makes the run a singleton: a second
concurrent run exits immediately rather than interleaving with the
first. The lock releases automatically on disconnect (session-scoped
locks do not survive the session), so no explicit unlock is needed.

\subsection{5b. TOC ingest pipeline}\label{b.-toc-ingest-pipeline}

\texttt{ingest\_toc(con)} follows the same shape as the metabase ingest,
with two differences. It calls \texttt{EUROSTATfetchR::fetch\_toc()}
(fetch + parse to flat \texttt{nodes} and \texttt{edges} frames), then:

\begin{enumerate}
\def\labelenumi{\arabic{enumi}.}
\tightlist
\item
  \textbf{Dedup guard} on \texttt{sha256} against the latest
  \texttt{toc\_snapshot} (rollback to an older tree is re-ingested, same
  as the metabase).
\item
  \textbf{Monotonicity guard:} refuse if the fetched
  \texttt{creationDate} is strictly older than the latest stored
  snapshot's (only possible because the TOC carries a timestamp).
\item
  \textbf{Diff and apply in one transaction, on two tables:}

  \begin{itemize}
  \tightlist
  \item
    \textbf{Nodes} --- a node is expired when its \texttt{code} is no
    longer present in staging \emph{with the same type and title}, and
    inserted when a staging row does not match a live row on all of
    \texttt{(code,\ type,\ title)}. So a title change expires the old
    row and inserts the new --- richer than the metabase's exact-tuple
    match, because the key (\texttt{code}) is narrower than the tracked
    columns. \texttt{title} comparison uses
    \texttt{IS\ NOT\ DISTINCT\ FROM} so null-title nodes round-trip
    correctly.
  \item
    \textbf{Edges} --- pure presence/absence on
    \texttt{(parent\_code,\ child\_code)}, exactly like the metabase
    pattern.
  \item
    A \textbf{bootstrap branch} (empty live set) bulk-inserts both
    tables with no anti-join, same rationale as the metabase.
  \end{itemize}
\end{enumerate}

The TOC ingest runs inside the main run's connection and advisory lock
(it does not take its own). It is folded into the scheduled job rather
than separately scheduled (Section 9), and runs \emph{before} the
metabase alert stage so folder resolution (Section 7.3) uses the
freshest tree. Its failure is isolated: a TOC ingest error emails the
maintainer but does \textbf{not} block metabase alerting --- folder
resolution then falls back to the last successfully ingested tree, which
is correct degradation.

\begin{center}\rule{0.5\linewidth}{0.5pt}\end{center}

\subsection{6. Change detection}\label{change-detection}

The change-detection read ---
\texttt{UMARaccessR::sql\_get\_eurostat\_metabase\_changes\_from\_snapshot(con,\ snapshot\_id)},
wrapping the \texttt{eurostat.get\_metabase\_changes} Postgres function
--- returns one row per affected dataset for the given snapshot. It
classifies at the dataset level first, then breaks down dimension and
level changes for datasets that were neither newly added nor wholly
removed. (A companion read,
\texttt{sql\_get\_eurostat\_removed\_levels\_from\_snapshot}, returns
the specific removed level codes for a set of datasets, used to list ---
not merely count --- removed levels in alert emails.)

\subsubsection{6.1 Event types}\label{event-types}

{\def\LTcaptype{none} % do not increment counter
\begin{longtable}[]{@{}
  >{\raggedright\arraybackslash}p{(\linewidth - 2\tabcolsep) * \real{0.5000}}
  >{\raggedright\arraybackslash}p{(\linewidth - 2\tabcolsep) * \real{0.5000}}@{}}
\toprule\noalign{}
\begin{minipage}[b]{\linewidth}\raggedright
\texttt{event}
\end{minipage} & \begin{minipage}[b]{\linewidth}\raggedright
Meaning
\end{minipage} \\
\midrule\noalign{}
\endhead
\bottomrule\noalign{}
\endlastfoot
\texttt{dataset\_added} & Dataset appeared for the first time (no prior
rows at all) \\
\texttt{dataset\_removed} & Dataset had live rows before, none survive
after \\
\texttt{changed} & Existing dataset with dimension and/or level
changes \\
\end{longtable}
}

For \texttt{changed} rows the query also returns:

\begin{itemize}
\tightlist
\item
  \texttt{dim\_changes} --- a text list like
  \texttt{geo:dim\_added,\ unit:dim\_removed}, or NULL if no
  dimension-level change.
\item
  \texttt{level\_added} / \texttt{level\_removed} --- counts of
  code-list positions added/removed (excluding \texttt{time}).
\end{itemize}

New and removed datasets are reported as a \textbf{single} event each;
their per-dimension changes are suppressed, because a whole-dataset
event should not also emit a \texttt{dim\_added}/\texttt{dim\_removed}
for every one of its dimensions.

\subsubsection{6.2 What counts as
``breaking''}\label{what-counts-as-breaking}

This is a \textbf{subscription} concept, not a detection concept, so it
is defined in the application (R), next to the subscription logic, not
in the SQL. A change is breaking if it can break an existing query:

\begin{itemize}
\tightlist
\item
  \texttt{dataset\_removed}
\item
  any dimension change (\texttt{dim\_added} or \texttt{dim\_removed})
\item
  \texttt{level\_removed\ \textgreater{}\ 0}
\end{itemize}

Not breaking:

\begin{itemize}
\tightlist
\item
  \texttt{level\_added} only --- a new code rarely breaks anything that
  already runs
\item
  \texttt{dataset\_added} --- a new dataset cannot break an existing
  query
\end{itemize}

\texttt{dim\_added} is treated as breaking because an added dimension
changes granularity, which silently breaks downstream aggregations and
joins --- this is precisely the failure that originally motivated the
project.

\begin{center}\rule{0.5\linewidth}{0.5pt}\end{center}

\subsection{7. Subscription system}\label{subscription-system}

\subsubsection{7.1 The spreadsheet}\label{the-spreadsheet}

Subscriptions are maintained by users themselves in a shared Excel file
on the network share.

The location is:
\texttt{"O:\textbackslash{}Avtomatizacija\textbackslash{}umar-automation-scripts\textbackslash{}data\textbackslash{}umar\_data\_bot\_subscriptions\textbackslash{}}
\texttt{eurostat\_email\_prijave.xlsx"}

Columns: \texttt{email}, \texttt{dataset}, \texttt{scope}, and an
optional free-text \texttt{notes} column that the system ignores
entirely. One row per \texttt{(email,\ dataset)}. Users add the dataset
codes they depend on; \texttt{scope} defaults to \texttt{breaking} if
left blank; \texttt{dataset\ =\ *} opts into everything.

Putting the sheet where the knowledge lives (with the people who know
their own dependencies) is deliberate --- it removes the maintainer as a
bottleneck. The spreadsheet is the human interface; the
\texttt{subscription} table is the validated, normalised truth derived
from it.

\subsubsection{7.2 Load and validation}\label{load-and-validation}

\texttt{load\_subscriptions(con)} reads the sheet and reloads the table.
It:

\begin{enumerate}
\def\labelenumi{\arabic{enumi}.}
\tightlist
\item
  \textbf{Selects columns by name}, so an extra column (like
  \texttt{notes}) cannot shift anything; a missing required column is a
  hard error.
\item
  \textbf{Normalises}: lowercases and trims email and dataset; defaults
  blank scope to \texttt{breaking}.
\item
  \textbf{Validates}, collecting problems without halting:

  \begin{itemize}
  \tightlist
  \item
    unknown codes --- each non-\texttt{*} code is checked against
    \emph{both} the live metabase datasets \emph{and} the live TOC
    folders; it is rejected only if it matches neither. This is the
    critical check; a typo'd code that matches nothing would otherwise
    silently deliver no alerts. Folder codes are entered exactly as they
    appear in the online Detailed-datasets tree (unprefixed); a
    \texttt{t\_}-prefixed Tables-hierarchy code will not match and is
    reported (the two hierarchies are genuinely different --- the system
    does not silently remap one to the other).
  \item
    invalid scope values,
  \item
    duplicate \texttt{(email,\ dataset)} pairs. Bad rows are
    \textbf{skipped, not fatal} --- one person's mistake must not block
    everyone else's subscriptions. Problems are emailed to the
    maintainer. Each surviving row is stamped with a derived
    \texttt{kind} (\texttt{dataset} or \texttt{folder}) from which
    lookup it matched.
  \end{itemize}
\item
  \textbf{Safety gate against corruption:} if the load would
  \emph{remove} existing subscriptions for \textbf{more than one email
  address}, it refuses (raises) rather than applying the change. The
  editing convention is that no one touches another person's rows, so
  removals spanning multiple people mean the sheet was mangled (cleared,
  truncated, saved empty). A single person editing their own rows passes
  fine. This protects the authoritative table from a corrupted sheet
  without a merge engine --- it is a sanity gate, not a diff.
\item
  \textbf{Truncate + reload} the clean rows, in one transaction. The
  sheet is authoritative, so idempotent full replacement is the simplest
  correct model; the gate handles the ``authoritative until someone
  breaks it'' case.
\end{enumerate}

\texttt{load\_subscriptions} runs on \textbf{every} run, not only on
change days, so that a corrupted sheet surfaces (via the gate) the same
night it happens rather than lying in wait until the next structural
change.

\subsubsection{7.3 Folder subscriptions}\label{folder-subscriptions}

A subscription code may be a \textbf{folder} rather than a dataset. A
folder subscription matches structural changes on \emph{any} dataset
beneath that folder in the TOC tree, so a user can subscribe to, say,
\texttt{mar} (maritime transport) and receive alerts for every
\texttt{mar\_*} dataset without listing them.

Resolution happens \textbf{at alert time, not load time}: the
subscription table stores the folder code as the user entered it, and
\texttt{filter\_alerts} resolves it to its descendant dataset/table
codes each run via
\texttt{UMARaccessR::sql\_resolve\_eurostat\_folder\_datasets} (a
recursive CTE over the \emph{live} TOC edges). The consequence --- and
the point --- is that a folder subscription \textbf{automatically covers
datasets added to that folder later}, with no re-subscription. Storing
the folder code (not a frozen expansion) also keeps the subscription
table honest: it records what the user asked for, and the current tree
is consulted at use.

The design keeps \textbf{dataset-code subscriptions as the reliable
backbone and folder subscriptions as an advisory layer on top.} A direct
dataset-code match never depends on the TOC; only folder matches do. So
if the TOC is momentarily stale (ingest failed, or Eurostat's tree lags
the metabase), a folder match may be slightly delayed but a direct code
match is unaffected --- and folder resolution can only ever point at
monitored datasets (verified: every dataset any folder resolves to is
present in the metabase).

\subsubsection{7.4 Automatic subscription of ingested pipeline
dependencies}\label{automatic-subscription-of-ingested-pipeline-dependencies}

\texttt{subscribe\_bot\_to\_ingested(con)} subscribes a dedicated bot
address to every Eurostat dataset the agency currently ingests, so a
structural change to any pipeline dependency is alerted on without
anyone having to maintain the list by hand. The ingested set is read
from \texttt{platform.table\ WHERE\ source\_id\ =\ 7} (Eurostat), which
is the source of truth for ``what we ingest''.

It runs \textbf{after} \texttt{load\_subscriptions} in \texttt{main},
because \texttt{load\_subscriptions} truncate-reloads the whole table
from the Excel sheet; the bot's rows come from \texttt{platform.table},
not the sheet, so they must be re-applied afterwards. It replaces only
its own rows (\texttt{DELETE\ WHERE\ email\ =\ bot\_email}, then
insert), so it never touches human subscriptions.

Two design points:

\begin{itemize}
\tightlist
\item
  \textbf{The bot address must be distinct from any human subscriber's},
  because the function owns and replaces all rows for that address every
  run. The bot uses a dot-free variant of the maintainer's Gmail
  address: Gmail ignores dots in the local part, so mail still lands in
  the same inbox, but the string differs from the dotted address a human
  would enter in the sheet --- so the two never collide. (A
  \texttt{+alias} would be more self-documenting; the dot-free form is
  functionally equivalent.)
\item
  \textbf{A metabase cross-check is the real value.} Every run, the
  function compares the ingested codes against the live metabase and
  emails the maintainer about any ingested table that is \textbf{not} in
  the current metabase. That is precisely a pipeline dependency that may
  have been renamed or removed upstream --- the \texttt{agr\_r\_animal}
  failure mode, where a dependency vanished and was only discovered
  weeks later by a failing test. Unlike a \texttt{dataset\_removed}
  alert (which is event-based --- it fires only on a live→gone
  transition the monitor observed, and only if subscribed), the
  cross-check is state-based: it catches a dependency that is
  \emph{already absent}, whenever and however it went, and does not
  depend on subscription state. It is a standing audit of the pipeline's
  foundations.
\end{itemize}

\begin{center}\rule{0.5\linewidth}{0.5pt}\end{center}

\subsection{8. Alerting}\label{alerting}

The orchestrator (\texttt{main}) runs, in order:

\begin{enumerate}
\def\labelenumi{\arabic{enumi}.}
\tightlist
\item
  Connect, take the advisory lock, configure email auth.
\item
  \texttt{load\_subscriptions} --- validate the human input first, fail
  fast if corrupt.
\item
  \texttt{subscribe\_bot\_to\_ingested} --- re-apply the bot's
  pipeline-dependency subscriptions and run the metabase cross-check
  (Section 7.4).
\item
  \texttt{ingest\_toc} --- refresh the TOC tree so folder resolution is
  current. Wrapped in its own \texttt{tryCatch} that emails the
  maintainer on failure but does \textbf{not} re-raise: a TOC failure
  must not block metabase alerting.
\item
  \texttt{ingest} (metabase) --- wrapped so any failure emails the
  maintainer and still records the traceback in the log.
\item
  If \texttt{unchanged}, exit.
\item
  If \texttt{changed}: read changes via
  \texttt{UMARaccessR::sql\_get\_eurostat\_metabase\_changes\_from\_snapshot}
  → filter per subscriber → send one email per affected subscriber. This
  whole stage is wrapped so a failure in the alert query or send emails
  the maintainer rather than dying silently in the log.
\end{enumerate}

\subsubsection{8.1 Per-person filtering}\label{per-person-filtering}

\texttt{filter\_alerts} joins the run's change rows to the subscription
table, resolving each distinct folder once per run and reusing the
result across subscribers:

\begin{itemize}
\tightlist
\item
  \texttt{dataset\ =\ \textquotesingle{}*\textquotesingle{}} matches all
  change rows.
\item
  \texttt{kind\ =\ \textquotesingle{}folder\textquotesingle{}} matches
  rows whose dataset is in the folder's resolved descendant set (Section
  7.3).
\item
  \texttt{kind\ =\ \textquotesingle{}dataset\textquotesingle{}} matches
  by exact dataset code.
\item
  \texttt{scope\ =\ \textquotesingle{}breaking\textquotesingle{}} then
  keeps only breaking rows (6.2);
  \texttt{scope\ =\ \textquotesingle{}all\textquotesingle{}} keeps
  everything for the matched datasets.
\end{itemize}

One email is sent per affected subscriber, containing only their matched
rows, and only if they have at least one; silence on non-events.

\subsubsection{8.2 Email format}\label{email-format}

Alerts are HTML email with the change table wrapped in
\texttt{\textless{}pre\textgreater{}} for a monospace font, so the
space-padded columns actually align.

A \texttt{dataset\_removed} row carries an added note suggesting the
dataset may have been renamed and to check for a replacement code. This
is a cheap mitigation of the rename-detection gap (Section 10): the
system cannot yet link a removal to its replacement, but it can tell the
affected subscriber to go look.

\subsubsection{8.3 Two distinct email
concerns}\label{two-distinct-email-concerns}

\begin{itemize}
\tightlist
\item
  \textbf{Operational failure} (ingest broke, sheet corrupt, alert query
  failed, pipeline dependency missing from metabase) → goes to the
  maintainer. Distinctive subject prefix for inbox filtering.
\item
  \textbf{Structural alert} (something changed) → goes to subscribers.
\end{itemize}

These are separate functions with separate recipients and must not be
merged.

\begin{center}\rule{0.5\linewidth}{0.5pt}\end{center}

\subsection{9. Operations}\label{operations}

\begin{itemize}
\tightlist
\item
  \textbf{Schedule:} Windows Task Scheduler runs the batch script
  \texttt{O:\textbackslash{}Avtomatizacija\textbackslash{}umar-automation-scripts\textbackslash{}batch-scripts\textbackslash{}eurostat-update-metabase.bat}
  a few minutes after the Eurostat evening refresh (\textasciitilde23:00
  CET), and after the 11:00 refresh. The task name is
  \texttt{eurostat-update-metabase}. Exact timing is not critical: the
  \texttt{sha256} dedup makes a catch-up run after downtime free (it
  exits immediately if nothing changed). The \textbf{TOC ingest is
  folded into this same job} rather than separately scheduled: the TOC
  regenerates \textasciitilde00:11 CET, but folder resolution only needs
  a reasonably current tree, so refreshing it twice daily alongside the
  metabase run is ample. (If TOC \emph{change alerts} --- folder
  renames, dataset moves --- are built later, revisit whether they
  warrant their own cadence.)
\item
  \textbf{Batch wrapper:} the batch script calls \texttt{Rscript} on the
  runner (\texttt{R\textbackslash{}eurostat-update-metabase.R}), which
  sources the function files and calls \texttt{main()} inside a
  \texttt{tryCatch}. Output is redirected to the log; the batch then
  \texttt{pushd}es into the repo and commits the log.
\item
  \textbf{Logging:} stdout/stderr are redirected to
  \texttt{logs\textbackslash{}eurostat\_metabase.Rout} on the share,
  which is committed to git for an audit trail. The \texttt{.Rout} is
  overwritten each run; git preserves the history. The log is \emph{not}
  the alert channel --- alerts and failures are emailed from within the
  script via \texttt{gmailr}.
\item
  \textbf{Alert channel is email only.} Failure surfaces as a maintainer
  email plus a log entry; there is no second channel.
\item
  \textbf{Secrets:} the DB password is read from an environment
  variable; the Gmail OAuth credentials and token cache are at absolute
  paths on the share. The scheduled account must be able to read the
  credentials and write the token cache.
\end{itemize}

\begin{center}\rule{0.5\linewidth}{0.5pt}\end{center}

\subsection{10. Known limitations and
backlog}\label{known-limitations-and-backlog}

Recorded honestly so a future maintainer knows the edges.

\begin{itemize}
\item
  \textbf{Dataset renames are detected but not yet linked (now enabled
  by the TOC).} A renamed dataset appears as an unrelated
  \texttt{dataset\_removed} (old code) and \texttt{dataset\_added} (new
  code); the metabase alone cannot connect them (codes, no titles). The
  \texttt{dataset\_removed} alert includes a ``check for a replacement''
  hint as a stopgap. The TOC is now ingested and carries titles, so the
  \emph{mechanism} for the fix now exists --- match a removed node's
  title to an added node's title to identify a rename near-decisively.
  \emph{Backlog:} build TOC change detection (node title changes =
  renames; edge changes = folder moves) and surface renamed datasets in
  the \texttt{dataset\_removed} alert. The infrastructure is in place;
  the detection query and wiring are not yet built.
\item
  \textbf{Level (code) renames are not distinguished from drop+add.} A
  code renamed within a dimension shows as \texttt{level\_removed} +
  \texttt{level\_added}, unlinked, for the same reason (no labels in the
  metabase).
\item
  \textbf{Dimension renames rely on a heuristic.} A dimension rename
  shows as \texttt{dim\_removed} + \texttt{dim\_added} in the same
  dataset and run; a matching ordered position signature makes it
  near-decisive, but it remains a candidate, not a certainty. (The
  signature logic was prototyped but is not currently wired into the
  change-detection function; fold it back in when a real dimension
  rename first occurs.)
\item
  \textbf{Reorder-without-rename is not detected.} If Eurostat reorders
  a code list without changing membership, no rows change, so nothing
  fires; \texttt{ord} on surviving rows also goes stale. Out of scope by
  choice --- reordering does not break code-based queries.
\item
  \textbf{The observation timestamp is fetch time, not publication time}
  for the metabase, because the endpoint provides no
  \texttt{Last-Modified}. Change \emph{ordering} is preserved by
  snapshot id; only the absolute ``when did Eurostat publish this'' is
  approximate. (The TOC, by contrast, carries a real
  \texttt{creationDate}.)
\item
  \textbf{OAuth token expiry is the most likely silent-death mode.} If
  the Gmail refresh token expires or the token cache becomes unwritable
  by the scheduled account, all alerting fails silently --- no alerts,
  no failure emails. \emph{Mitigation:} periodically confirm alerts
  still arrive; consider a heartbeat or second channel if the system
  becomes critical.
\item
  \textbf{Comext data monitoring gap.} Comext (detailed trade,
  \texttt{ds-*}) is a separate Eurostat dissemination system, present in
  \emph{neither} the metabase nor the TOC. Structural changes to those
  datasets are not detected, and they cannot be folder-subscribed. They
  are also not currently ingested by EUROSTATfetchR, so it is a moot
  point for now. If Comext data becomes a dependency, monitoring it
  would mean a third parallel SCD2 source pointed at the Comext
  catalogue.
\item
  \textbf{The TOC XML omits some tables that exist in the metabase.} A
  handful of tables (e.g.~\texttt{tet00071}) appear in the metabase and
  in the online browser but are simply absent from the \texttt{toc/xml}
  file --- confirmed by grepping the raw source, not a parser fault (the
  browser stitches in listings the machine- readable file does not
  contain, the same way it does for Comext). Consequence: these tables
  are \textbf{folder-unreachable but code-reachable} --- a folder
  subscription will not surface them, but a direct dataset-code
  subscription will (they are monitored in the metabase). This is a
  Eurostat source- completeness gap, not fixable on our side, and
  benign: the affected codes are obscure \texttt{t*} tables, and the
  dataset-code backbone covers them.
\item
  \textbf{TOC change detection is not built.} The TOC is snapshotted
  with full SCD2 history, but nothing yet \emph{diffs} consecutive TOC
  snapshots to alert on folder renames, dataset moves between folders,
  or (per the item above) dataset renames. Folder subscriptions use the
  TOC only for \emph{resolution}, always against the current tree.
  Building TOC change alerts is backlog; the history is being
  accumulated (including the no-op snapshots noted in 4.5) so it will be
  available when built.
\end{itemize}

\begin{center}\rule{0.5\linewidth}{0.5pt}\end{center}

\subsection{11. Design principles applied
throughout}\label{design-principles-applied-throughout}

For future changes, the principles this system was built on:

\begin{itemize}
\tightlist
\item
  \textbf{Detect, don't schedule.} Eurostat publishes no change
  schedule; monitoring is necessarily proactive and
  after-the-fact-but-before-consumption.
\item
  \textbf{Fail loud.} Every stage that can fail either raises or emails
  the maintainer; nothing swallows an error silently. The one thing a
  monitor must never do is fail quietly.
\item
  \textbf{State audits complement event alerts.} Event-based detection
  (a change observed between two snapshots) misses anything already
  broken before the monitor looked or before a subscription existed; the
  ingested-dependency cross-check (7.4) is a state audit that catches
  those. Both are needed.
\item
  \textbf{One source of truth per fact.} Redundant columns encoding the
  same thing (a flag plus a timestamp; a stored classification plus the
  data it derives from) were rejected in favour of deriving one from the
  other. Derived-and-cached is acceptable; stored-as-a-second-primary
  drifts.
\item
  \textbf{Content identity over metadata.} Dedup on the file's hash, not
  on unreliable timestamps or headers.
\item
  \textbf{Keys over coincidences.} Lineage joins on integer snapshot
  ids, never on timestamp equality.
\item
  \textbf{Complexity must be justified by scale.} At \textasciitilde1.3M
  rows, several performance defences (fingerprint tiers, full-snapshot
  storage avoidance) were considered and dropped as unnecessary. The
  design is sized to the actual data, not an imagined order of magnitude
  larger.
\item
  \textbf{Reusable components are packages with tests; applications are
  scripts.} The DB read functions (UMARaccessR) and the Eurostat
  fetch/parse (EUROSTATfetchR) carry roxygen, SQL tests, and
  dittodb/fixture tests; the ingest orchestration, filtering, and email
  are application glue and stay in a standalone script. The test that a
  thing belongs in a package is genuine reuse --- the metabase and TOC
  ingests, though substantial, are source-specific jobs and stayed in
  the app.
\item
  \textbf{Store what the user said; resolve at use.} Folder
  subscriptions store the folder code the user entered and resolve it
  against the current tree at alert time, rather than freezing an
  expansion at load time --- so newly-added datasets are covered
  automatically and the stored record stays honest.
\end{itemize}

\end{document}
